\section{Teste de aceitação do preço}
\label{sec:teste}

\subsection{Objetivo e hipóteses}

Verificar se o preço da assinatura e o modelo de comissão são aceitos pelos públicos
pagantes. As hipóteses a testar são:

\begin{description}[font=\bfseries, leftmargin=1.4em, style=unboxed]
  \item[H1.] Pelo menos 60\% dos tutores entrevistados consideram o PetCuida+ a R\$~12,90/mês
  ``caro, mas aceitável'' ou ``barato'' depois de conhecerem os benefícios.
  \item[H2.] Pelo menos 20\% dos tutores declaram intenção de assinar a esse preço.
  \item[H3.] Pelo menos 2 de 3 clínicas abordadas aceitam testar o Plano Parceiro após o período gratuito.
\end{description}

\subsection{Método proposto}

\begin{enumerate}
  \item \textbf{Pesquisa rápida com tutores} (formulário online, 3 a 5 minutos), com amostra-alvo
  de 30 a 50 respostas. Os respondentes vêm dos oito entrevistados da Atividade~6, de ONGs
  e protetores parceiros e de grupos de bairro. A pesquisa combina:
  \begin{itemize}
    \item perguntas de sensibilidade a preço (Van Westendorp): em que valor mensal o serviço seria
    barato demais, uma pechincha, caro e caro demais;
    \item intenção de assinar a R\$~9,90, R\$~12,90 e R\$~14,90 (Gabor-Granger simplificado);
    \item opinião sobre a comissão de 10\% e sobre a tabela social.
  \end{itemize}
  \item \textbf{Perguntas de preço ao final do teste de uso}, em 5 a 8 sessões curtas com o
  aplicativo demonstrativo, com o roteiro de usabilidade da Atividade~6.
  \item \textbf{Conversa com clínicas} (3 parceiros): apresentação do Plano Parceiro de R\$~79,90/mês
  com 3 meses gratuitos e pergunta sobre a disposição de pagar após o período.
\end{enumerate}

\subsection{Critérios de decisão}

\begin{longtable}{L{6.0cm} L{8.2cm}}
\toprule
\cabfundo \cab{Resultado do teste} & \cab{Decisão} \\
\midrule
\endfirsthead
\toprule
\cabfundo \cab{Resultado do teste} & \cab{Decisão} \\
\midrule
\endhead
\bottomrule
\endfoot
H1 e H2 confirmadas & Manter R\$~12,90/mês e R\$~119,00/ano \\
H1 confirmada, H2 não & Manter o preço e melhorar a proposta do PetCuida+ (benefícios e comunicação) \\
Preço aceitável abaixo de R\$~12,90 & Reduzir para R\$~9,90 e rever o ponto de equilíbrio (Tabela~\ref{tab:sens}) \\
H3 não confirmada & Estender o período gratuito ou reduzir a mensalidade do Plano Parceiro \\
\end{longtable}

\subsection{Cuidados}

Os respondentes devem concordar com o uso de suas respostas (LGPD), e a pesquisa deve evitar
perguntas indutoras. As respostas e a data de coleta ficam registradas para consulta futura.

\subsection{Roteiro resumido do questionário}

\begin{enumerate}
  \item Você tem cão ou gato? Quantos? Em qual bairro mora?
  \item Quanto gastou com o pet nos últimos 12 meses (aproximadamente)?
  \item Já deixou de levar o pet ao veterinário por falta de dinheiro? (sim/não)
  \item Descrição do PetCuida+ (lembretes, prontuário, descontos, créditos solidários).
  \item A que valor mensal você acharia o serviço barato demais? Uma pechincha? Caro? Caro demais?
  \item Você assinaria a R\$~9,90? E a R\$~12,90? E a R\$~14,90? (sim/talvez/não)
  \item O que faria você pagar mais ou menos?
\end{enumerate}
